\chapter{Introduction}
\pagenumbering{arabic}

The Internet of Things (IoT) has expanded rapidly and now underpins critical
applications across smart manufacturing, power grids, remote health monitoring,
and intelligent transport systems.
These deployments rely on large numbers of resource-constrained devices communicating
over shared wireless channels that are inherently exposed to eavesdropping.
Securing such networks against impersonation, replay, man-in-the-middle, key
compromise, and device-capture attacks is therefore a fundamental operational
requirement.

Considerable research has addressed IoT device authentication through mechanisms
tailored to the severe power and memory constraints of embedded hardware.
Physical Unclonable Functions (PUFs), hash-based challenge--response protocols, and
symmetric-key constructions have emerged as dominant lightweight approaches owing to
their low computational overhead.
DAuth~\cite{das2026comsnets} follows this direction: it combines PUFs with a
hash-based accumulator, eliminates long-term key storage, distributes authentication
across multiple in-network nodes, and supports multihop verification at minimal
overhead.

Despite these properties, a fundamental privacy gap persists.
During each authentication session, the device's real identity $ID_D$ is transmitted
in plaintext over the public channel.
A passive adversary can capture these transmissions to track the device across
sessions, correlate repeated authentication attempts, and infer network
behaviour, which directly violates modern privacy requirements.

Identity privacy has been established as a fundamental property in IoT
authentication, not merely an optional feature.
Survey literature confirms that anonymity, untraceability, and forward secrecy are
essential, particularly when devices operate in open or adversarial
environments~\cite{zhong2025survey}.
Existing countermeasures such as session-refreshed pseudonyms and hash-chain-based
temporary identities partially address this exposure but often introduce new
challenges such as desynchronisation, excessive storage overhead, slow identity
lookup, or increased computation.
The sub-framework taxonomy of~\cite{zhong2025survey}, covering SF3-UA through
SF9-FS, illustrates how difficult it is to simultaneously satisfy anonymity,
unlinkability, and resilience to message loss.

This work investigates how anonymity can be integrated into DAuth without
compromising its lightweight properties.
All authentication messages of DAuth traverse an open channel, and transmitting
$ID_D$ in plaintext exposes the device to tracking, session linkage, and targeted
impersonation.
Drawing on prior anonymous-authentication literature, this work identifies a
low-overhead pseudonym mechanism that integrates cleanly into DAuth.
The objective is identity anonymity, session unlinkability, and non-traceability,
achieved without relaxing the lightweight and distributed properties of DAuth.

A second gap in DAuth concerns desynchronisation resilience.
Because DAuth retains only a single session nonce at the authentication server,
any loss of the key-exchange response packet causes the device and server states to
diverge, leading to a permanent authentication failure that can only be resolved
through costly full re-enrolment.

\section{Contributions}

This work makes the following key contributions.
\begin{enumerate}
  \item An anonymity-preserving scheme is proposed that introduces per-session
    pseudonym rotation to preserve device anonymity and a dual-state mechanism
    to recover from desynchronisation without re-enrolment.
  \item The security of the proposed scheme is validated through both informal
    analysis and formal verification using ProVerif.
  \item The efficiency of the proposed scheme is evaluated using theoretical
    analysis, COOJA\slash Contiki-NG simulation, and hardware experiments on
    Raspberry Pi~4B, covering communication overhead, energy consumption, and
    CPU utilisation.
\end{enumerate}

\section{Organisation of The Report}

The rest of this report is organised as follows.
Chapter~2 surveys the related literature.
Chapter~3 describes the prerequisite concepts.
Chapter~4 overviews DAuth~\cite{das2026comsnets}.
Chapter~5 presents the proposed scheme in full detail.
Chapter~6 contains the security analysis.
Chapter~7 presents the performance evaluation.
Chapter~8 concludes the report.
